% ************************************************************
% LITERATURE REVIEW
% ************************************************************
\section{Literature Review}

\subsection{WASH Disruption as a Public-Health Systems Problem}

The WASH literature treats water, sanitation, and hygiene as linked determinants of infectious-disease risk, health-system resilience, dignity, and equity. Reviews of WASH interventions have consistently connected safe water, sanitation, and hygiene to diarrhoeal-disease prevention, while also showing that intervention effects vary by context, implementation quality, baseline infrastructure, and population vulnerability \citep{fewtrell_2005,cairncross_2010,taylor_2015}. In humanitarian settings, this context dependence is central because exposure is shaped not only by household behavior but also by damaged infrastructure, crowding, service interruption, fuel availability, and restricted mobility. A Gaza WASH--health priority index is therefore most defensible when framed as a systems-level prioritization tool rather than as a direct estimate of individual disease incidence.

Humanitarian WASH standards also support an integrated approach. The Sphere Handbook links emergency WASH response to water supply, excreta management, hygiene promotion, vector control, solid-waste management, and outbreak risk reduction \citep{sphere_2018}. This integrated framing is important for Gaza because public reporting documents simultaneous stress across drinking water access, sewage exposure, solid waste, pests, displacement-site crowding, and health-service constraints \citep{ocha_2026_apr17,ocha_2026_may15}. A single-domain indicator, such as water delivery alone, would understate the public-health risk created by the convergence of multiple environmental and service-system failures.

\subsection{Humanitarian Settings, Displacement, and Environmental Exposure}

Conflict and displacement settings create conditions that intensify communicable-disease risk: population movement, crowded shelters, disrupted primary care, damaged water systems, reduced sanitation, solid-waste accumulation, and constrained surveillance. The complex-emergencies literature emphasizes that communicable-disease control depends on basic public-health infrastructure, functioning health services, surveillance, and coordinated response capacity \citep{connolly_2004}. This is directly relevant to Gaza, where displacement-site indicators include environmental hazards and reported skin infections or infestations, while broader humanitarian reporting documents limits on water, sanitation, fuel, maintenance, and service continuity \citep{ocha_2026_apr17,ocha_2026_may15}.

The literature also supports caution in interpretation. WASH indicators in crisis settings often represent exposure signals rather than confirmed clinical outcomes. Site-level reports of sewage, pests, standing water, or solid waste can identify plausible environmental-health risk but cannot, by themselves, establish individual-level incidence or causal effects. For that reason, a rapid index should use conservative language: risk convergence, prioritization, service stress, and public-health signal. This approach preserves analytic utility while avoiding unsupported claims about attributable disease burden.

\subsection{Spatial Prioritization and Secondary Data}

Humanitarian public-health decision-making often depends on incomplete, rapidly changing secondary data. In that setting, the analytic requirement is not perfect measurement, but transparent source handling, explicit geography, reproducible extraction, and conservative interpretation. Public data can be organized into an auditable prioritization framework when source provenance, extraction rules, scoring assumptions, and verification status are preserved. This logic supports the present staged design: first, build a transparent index from accessible public reports; second, validate and refine it with facility, boundary, and WASH-service datasets as they become exportable and verifiable.

Gaza is a suitable case for this approach because the first-pass public data are uneven but actionable. OCHA reports provide date-stamped evidence of displacement-site hazards, WASH facility generator dependence, sewage flooding in Khan Younis, and operational constraints \citep{ocha_2026_apr17,ocha_2026_may15}. WHO, WASH Cluster, and HDX sources provide candidate validation layers for health-service functionality, WASH response indicators, and geospatial denominators \citep{who_emro_dashboards_2026,wash_cluster_palestine_2026,hdx_palestine_health_2026}. The literature therefore supports a staged design that starts with accessible public reports and then adds more precise spatial and facility data once source fields can be verified.

\subsection{Gap Addressed by the Present Study}

The current gap is not the absence of descriptive reporting on Gaza. The gap is the lack of a compact, reproducible method for translating fragmented public reporting into a transparent public-health prioritization structure. Existing humanitarian reports document hazards, service constraints, and infrastructure stress, but these indicators are often dispersed across situation reports, dashboards, and geospatial portals. The present study addresses this gap by formalizing a five-domain WASH--health priority index with explicit source eligibility rules, manual extraction fields, domain scoring, and sensitivity analysis. This design is appropriate for early-stage public-health intelligence work because it produces an auditable structure that can be corrected, expanded, or replaced as better data become available.
% ************************************************************
% END LITERATURE REVIEW
% ************************************************************
